%----------------------------------------------------------
% 제11장. AI 거버넌스 감사 및 인증
%----------------------------------------------------------

\chapter{AI 거버넌스 감사 및 인증}
\label{ch:audit_certification}

AI 거버넌스 체계가 구축되었다면, 
그것이 실제로 작동하고 있는지 검증하는 것이 다음 단계이다. 
감사(감사)는 거버넌스의 실효성을 확인하고, 
개선점을 발굴하며, 이해관계자에게 신뢰를 제공\cite{toreini2020relationship}하는 핵심 메커니즘이다.

본 장에서는 AI 거버넌스 감사의 유형, 방법론, 수행 절차를 다루고, 인증 제도에 대한 준비 방법\cite{zabala2024state}을 제시한다. 특히 알고리즘의 책임성 공백을 메우기 위한 통합 감사 프레임워크\cite{raji2020closing}와 실무적 감사 기법\cite{raji2019actionable}을 중심으로 실무 지침을 구성한다.


%----------------------------------------------------------
\section{AI 거버넌스 감사 개요}
\label{sec:11.1}

\subsection{감사의 목적과 가치}
\label{sec:11.1.1}

AI 거버넌스 감사는 다음과 같은 목적을 가진다.

\begin{enumerate}
 \item \textbf{컴플라이언스 확인\cite{mokander2022conformity}}: 
 내부 정책, 외부 규제, 산업 표준에 대한 준수 여부 검증
 
 \item \textbf{리스크 식별}:
 거버넌스 체계의 취약점, 갭, 잠재적 리스크 발굴
 
 \item \textbf{효과성 평가}: 
 거버넌스 활동이 실제로 의도한 효과를 달성하는지 평가
 
 \item \textbf{개선 기회 도출}: 
 현재 수준에서 더 나은 수준으로 발전하기 위한 방향 제시
 
 \item \textbf{신뢰 제공}: 
 경영진, 이사회, 규제 기관, 고객에게 거버넌스 신뢰성 입증
\end{enumerate}


\subsection{감사 유형}
\label{sec:11.1.2}

\begin{table}[htbp]
\centering
\caption{AI 거버넌스 감사 유형}
\label{tab:audit_types}
\begin{tabularx}{\textwidth}{p{2.5cm}p{2cm}p{2.5cm}X}
\toprule
\textbf{유형} & \textbf{수행 주체} & \textbf{주기} & \textbf{특징} \\
\midrule
자가 Check & 
시스템 담당자 & 
분기별 & 
일상적 Check, 빠른 피드백, 부담 최소화 \\
\midrule
내부 감사 & 
내부감사팀 또는 거버넌스 조직 & 
연간 & 
독립적 시각, 전사 관점, 개선 권고 \\
\midrule
2차 감사 (Second-Party) & 
고객, 파트너 & 
계약 based & 
비즈니스 관계에서 요구, 신뢰 구축 \\
\midrule
외부 감사 (Third-Party) & 
독립 감사 기관 & 
2년 주기 또는 규제 요구 시 & 
최고 수준 객관성, 인증 연계 가능. 외부 독립 감사는 AI 시스템의 안전성을 객관적으로 입증하는 핵심 수단이다.\cite{falco2021governing} \\
\midrule
규제 감사 & 
규제 기관 & 
불시 또는 정기 & 
법적 준수 확인, 제재 가능 \\
\bottomrule
\end{tabularx}
\end{table}


\subsection{감사 Scope}
\label{sec:11.1.3}

AI 거버넌스 감사의 범위는 다음 영역을 포함한다.

\begin{table}[htbp]
\centering
\caption{AI 거버넌스 감사 Scope}
\label{tab:audit_scope}
\begin{tabularx}{\textwidth}{p{3cm}X}
\toprule
\textbf{영역} & \textbf{세부 항목} \\
\midrule
정책 및 전략 & 
AI 윤리 정책 존재 및 승인 여부; 정책의 현행화; 
전략과의 정합성 \\
\midrule
조직 및 역할 & 
거버넌스 조직 구성; 역할.책임 명확성; 
Resource 충분성; 독립성 \\
\midrule
프로세스 & 
개발-배포-운영 프로세스; 승인 Gate; 
변경 관리; Retire 절차 \\
\midrule
기술 제어 & 
공정성\cite{bellamy2019ai} 평가; 설명가능한 AI(설명가능한 AI(XAI)) 구현; 모니터링\cite{klaise2021monitoring}; 
보안; Privacy 보호 \\
\midrule
문서화 & 
모델 카드; 영향평가; 승인 기록; 
운영 로그; 5년 보관 준수 \\
\midrule
인적 감독 & 
감독 수준 적절성; 감독자 역량; 
번복 메커니즘 작동 \\
\midrule
이용자 권리 & 
설명 제공; 이의 제기 절차; 
구제 메커니즘 \\
\midrule
인시던트 관리 & 
Response 절차; 보고 체계; 
재발 방지 조치 \\
\bottomrule
\end{tabularx}
\end{table}


%----------------------------------------------------------
\section{내부 감사 수행}
\label{sec:11.2}

\subsection{감사 계획 수립}
\label{sec:11.2.1}

\begin{lstlisting}[language=Python, caption={AI Governance Audit Plan Generator}, label={lst:audit_plan_generator}]
"""
AI Governance Audit Plan Generator
"""

from dataclasses import dataclass, field
from typing import Dict, List, Optional
from datetime import date, timedelta
from enum import Enum

class AuditType(Enum):
 """Audit Types"""
 SELF_ASSESSMENT = "Self-Assessment"
 INTERNAL = "Internal Audit"
 EXTERNAL = "External Audit"
 REGULATORY = "Regulatory Audit"

class RiskLevel(Enum):
 """Risk Levels"""
 HIGH = "High Risk"
 MEDIUM = "Medium Risk"
 LOW = "Low Risk"

@dataclass
class AuditScope:
 """Audit Scope"""
 systems: List[str] # Target System
 domains: List[str] # Audit 
 period: str # Target Period

@dataclass
class AuditPlan:
 """Audit """
 audit_id: str
 audit_name: str
 audit_type: AuditType
 
 # days 
 planning_start: date
 fieldwork_start: date
 fieldwork_end: date
 report_date: date
 
 # Scope
 scope: AuditScope
 
 # 
 lead_auditor: str
 audit_team: List[str]
 
 # 
 criteria: List[str]
 
 # Risk based Priority
 priority_areas: List[Dict[str, str]] = field(default_factory=list)


class AuditPlanGenerator:
 """Audit Generator"""
 
 # Standard Audit Domains
 STANDARD_DOMAINS = [
 "Policy & Strategy",
 "Organization & Roles",
 "Development Process",
 "Deployment & Approval",
 "Operation & Monitoring",
 "Fairness Management",
 "Explainability",
 "Human Oversight",
 "Documentation & Records",
 "User Rights",
 "Incident Management",
 ]
 
 # Standard Audit Criteria
 STANDARD_CRITERIA = [
 "AI Basic Act & Enforcement Decree",
 "Internal AI Ethics Policy",
 "Internal AI Governance Standards",
 "Industry-specific Regulations (if applicable)",
 "ISO/IEC 42001\cite{iso42001} (Reference)",
 ]
 
 def __init__(self, organization_name: str):
 self.organization = organization_name
 
 def generate_annual_plan(
 self,
 year: int,
 ai_systems: List[Dict],
 ) -> List[AuditPlan]:
 """Generate annual audit plan"""
 
 plans = []
 
 # Identify high-risk systems
 high_risk = [s for s in ai_systems if s.get('risk_level') == 'high']
 medium_risk = [s for s in ai_systems if s.get('risk_level') == 'medium']
 low_risk = [s for s in ai_systems if s.get('risk_level') == 'low']
 
 # Q1: High-risk system audit
 if high_risk:
 plans.append(AuditPlan(
 audit_id=f"AUD-{year}-001",
 audit_name=f"{year} High-Risk AI System Audit",
 audit_type=AuditType.INTERNAL,
 planning_start=date(year, 1, 15),
 fieldwork_start=date(year, 2, 1),
 fieldwork_end=date(year, 2, 28),
 report_date=date(year, 3, 15),
 scope=AuditScope(
 systems=[s['system_id'] for s in high_risk],
 domains=self.STANDARD_DOMAINS,
 period=f"Entire {year-1}",
 ),
 lead_auditor="TBD",
 audit_team=[],
 criteria=self.STANDARD_CRITERIA,
 priority_areas=[
 {"area": "Fairness Management", "reason": "Mandatory for high-risk"},
 {"area": "Human Oversight", "reason": "Regulatory requirement"},
 {"area": "Documentation & Records", "reason": "5-year retention obligation"},
 ],
 ))
 
 # Q2: Governance framework audit
 plans.append(AuditPlan(
 audit_id=f"AUD-{year}-002",
 audit_name=f"{year} AI Governance Framework Audit",
 audit_type=AuditType.INTERNAL,
 planning_start=date(year, 4, 1),
 fieldwork_start=date(year, 4, 15),
 fieldwork_end=date(year, 5, 15),
 report_date=date(year, 5, 31),
 scope=AuditScope(
 systems=["Enterprise"],
 domains=["Policy & Strategy", "Organization & Roles", "Incident Management"],
 period=f"July {year-1} - March {year}",
 ),
 lead_auditor="TBD",
 audit_team=[],
 criteria=self.STANDARD_CRITERIA[:3],
 priority_areas=[
 {"area": "Policy Update", "reason": "Reflect regulatory changes"},
 {"area": "Role Implementation", "reason": "Verify effectiveness"},
 ],
 ))
 
 # Q3: Medium-risk system sample audit
 if medium_risk:
 sample_size = min(5, len(medium_risk))
 sample = medium_risk[:sample_size]
 
 plans.append(AuditPlan(
 audit_id=f"AUD-{year}-003",
 audit_name=f"{year} Medium-Risk AI System Sample Audit",
 audit_type=AuditType.INTERNAL,
 planning_start=date(year, 7, 1),
 fieldwork_start=date(year, 7, 15),
 fieldwork_end=date(year, 8, 15),
 report_date=date(year, 8, 31),
 scope=AuditScope(
 systems=[s['system_id'] for s in sample],
 domains=self.STANDARD_DOMAINS[:7],
 period=f"First half of {year}",
 ),
 lead_auditor="TBD",
 audit_team=[],
 criteria=self.STANDARD_CRITERIA[:3],
 priority_areas=[
 {"area": "Monitoring", "reason": "Drift management"},
 ],
 ))
 
 # Q4: Annual Comprehensive Assessment
 plans.append(AuditPlan(
 audit_id=f"AUD-{year}-004",
 audit_name=f"{year} AI Governance Comprehensive Assessment",
 audit_type=AuditType.INTERNAL,
 planning_start=date(year, 10, 1),
 fieldwork_start=date(year, 10, 15),
 fieldwork_end=date(year, 11, 15),
 report_date=date(year, 11, 30),
 scope=AuditScope(
 systems=["Enterprise"],
 domains=self.STANDARD_DOMAINS,
 period=f"Entire {year}",
 ),
 lead_auditor="TBD",
 audit_team=[],
 criteria=self.STANDARD_CRITERIA,
 priority_areas=[
 {"area": "Maturity Assessment", "reason": "Measure annual progress"},
 {"area": "Remediation", "reason": "Follow-up on previous findings"},
 ],
 ))
 
 return plans
 
 def generate_plan_document(self, plan: AuditPlan) -> str:
 """Generate audit plan document"""
 
 doc = f"""
# AI Governance Audit Plan

---

## 1. Audit 

| Item | Content |
|------|------|
| Audit ID | {plan.audit_id} |
| Audit Name | {plan.audit_name} |
| Audit Type | {plan.audit_type.value} |
| Lead Auditor | {plan.lead_auditor} |

---

## 2. Audit days 

| Phase | Start | End |
|------|------|------|
| Planning | {plan.planning_start} | {plan.fieldwork_start - timedelta(days=1)} |
| Fieldwork | {plan.fieldwork_start} | {plan.fieldwork_end} |
| Reporting | {plan.fieldwork_end + timedelta(days=1)} | {plan.report_date} |

---

## 3. Audit Scope

### 3.1 Target System

"""
 
 for system in plan.scope.systems:
 doc += f"- {system}\n"
 
 doc += f"""

### 3.2 Audit 

"""
 
 for domain in plan.scope.domains:
 doc += f"- {domain}\n"
 
 doc += f"""

### 3.3 Target Period

{plan.scope.period}

---

## 4. Audit 

"""
 
 for i, criterion in enumerate(plan.criteria, 1):
 doc += f"{i}. {criterion}\n"
 
 doc += f"""

---

## 5. Risk based Priority

| | Priority |
|------|--------------|
"""
 
 for area in plan.priority_areas:
 doc += f"| {area['area']} | {area['reason']} |\n"
 
 doc += f"""

---

## 6. Audit 

| Role | |
|------|--------|
| Audit | {plan.lead_auditor} |
"""
 
 for i, member in enumerate(plan.audit_team, 1):
 doc += f"| Audit {i} | {member} |\n"
 
 doc += """

---

## 7. 

1. Audit Report
2. 3. 4. ---

## 8. Approval

| Role | Name | Signature | Date |
|------|------|------|------|
| Lead Auditor | | | |
| Head of Audit | | | |
| Dept. Head | | | |

---

*This audit plan is subject to change.*
"""
 
 return doc
\end{lstlisting}


\subsection{감사 수행 방법론}
\label{sec:11.2.2}

AI 거버넌스 감사는 다음 단계로 수행된다.

\paragraph{1단계: 사전 준비}

\begin{itemize}
 \item 감사 범위 및 목표 확정
 \item 감사 기준(법규, 정책, 표준) 확인
 \item 사전 자료 요청 및 검토
 \item 감사 프로그램(질문지, 체크리스트) 준비
 \item 킥오프 미팅
\end{itemize}

\paragraph{2단계: 현장 감사}

\begin{itemize}
 \item 문서 검토: 정책, 절차, 기록 확인
 \item 인터뷰: 담당자, 개발자, 운영자 면담
 \item 시스템 검증: 실제 시스템 작동 확인
 \item 샘플 테스트: 데이터, 결정, 설명 샘플 검토
 \item 워크스루: 프로세스 전체 따라가기
\end{itemize}

\paragraph{3단계: 분석 및 평가}

\begin{itemize}
 \item 증거 분석 및 검증
 \item 기준 대비 준수 여부 결정
 \item 지적사항\cite{cobbe2021reviewable} 도출 및 등급 분류
 \item 근본 원인 분석
 \item 개선 권고사항 도출
\end{itemize}

\paragraph{4단계: 보고}

\begin{itemize}
 \item Draft 작성 및 피감사 부서 검토
 \item 이견 조율
 \item 최종 Report 작성
 \item 경영진 보고
\end{itemize}


\subsection{감사 체크리스트}
\label{sec:11.2.3}

\begin{lstlisting}[language=Python, caption={AI Governance Audit Checklist}, label={lst:audit_checklist}]
"""
AI Governance Audit Checklist

Detailed Check Items by Area
"""

from dataclasses import dataclass, field
from typing import Dict, List, Optional
from enum import Enum

class ComplianceStatus(Enum):
 """Compliance Status"""
 COMPLIANT = "Compliant"
 PARTIALLY_COMPLIANT = "Partially Compliant"
 NON_COMPLIANT = "Non-Compliant"
 NOT_APPLICABLE = "N/A"
 NOT_TESTED = "Not Tested"

@dataclass
class AuditCheckItem:
 """Audit Check Items"""
 item_id: str
 domain: str
 category: str
 question: str
 criteria: str
 
 # Check 
 status: ComplianceStatus = ComplianceStatus.NOT_TESTED
 evidence: str = ""
 finding: str = ""
 recommendation: str = ""


class AIGovernanceAuditChecklist:
 """AI Governance Audit Checklist"""
 
 def __init__(self):
 self.items: List[AuditCheckItem] = self._initialize_checklist()
 
 def _initialize_checklist(self) -> List[AuditCheckItem]:
 """ Checklist """
 
 items = []
 
 # 1. Policy & Strategy
 items.extend([
 AuditCheckItem(
 item_id="POL-01",
 domain="Policy & Strategy",
 category="Policy Existence",
 question="Is there a formal AI Ethics Policy?",
 criteria="Documented policy, executive approval",
 ),
 AuditCheckItem(
 item_id="POL-02",
 domain="Policy & Strategy",
 category="Policy Maintenance",
 question="Has the policy been reviewed/updated within the last year?",
 criteria="Annual review cycle",
 ),
 AuditCheckItem(
 item_id="POL-03",
 domain="Policy & Strategy",
 category="Policy Dissemination",
 question="Has the policy been disseminated to relevant staff?",
 criteria="Training records, awareness check",
 ),
 AuditCheckItem(
 item_id="POL-04",
 domain="Policy & Strategy",
 category="Regulatory Alignment",
 question="Does the policy reflect the latest regulations?",
 criteria="Tracking and reflecting legal changes",
 ),
 ])
 
 # 2. Role
 items.extend([
 AuditCheckItem(
 item_id="ORG-01",
 domain=" Role",
 category=" ",
 question="AI Governance / ?",
 criteria=" , ",
 ),
 AuditCheckItem(
 item_id="ORG-02",
 domain=" Role",
 category="Role Definition",
 question="RACI Role Definition ?",
 criteria="RACI Document, ",
 ),
 AuditCheckItem(
 item_id="ORG-03",
 domain=" Role",
 category="Resource ",
 question="Governance ?",
 criteria=" Analysis",
 ),
 AuditCheckItem(
 item_id="ORG-04",
 domain=" Role",
 category=" ",
 question=" ( , ) ?",
 criteria=" , ",
 ),
 ])
 
 # 3. AI System Management
 items.extend([
 AuditCheckItem(
 item_id="SYS-01",
 domain="AI System Management",
 category=" ",
 question=" AI System ?",
 criteria=" ",
 ),
 AuditCheckItem(
 item_id="SYS-02",
 domain="AI System Management",
 category=" Classification",
 question=" System Level Classification ?",
 criteria=" ",
 ),
 AuditCheckItem(
 item_id="SYS-03",
 domain="AI System Management",
 category="Model ",
 question=" System Model ?",
 criteria="Model ",
 ),
 ])
 
 # 4. items.extend([
 AuditCheckItem(
 item_id="DEV-01",
 domain=" ",
 category="Approval Gate",
 question=" Approval Operation ?",
 criteria="Approval ",
 ),
 AuditCheckItem(
 item_id="DEV-02",
 domain=" ",
 category="Approval Compliance",
 question=" Approval ?",
 criteria="Approval 0 ",
 ),
 AuditCheckItem(
 item_id="DEV-03",
 domain=" ",
 category="Data ",
 question="Learning Data Management ?",
 criteria="Data Check ",
 ),
 ])
 
 # 5. Fairness
 items.extend([
 AuditCheckItem(
 item_id="FAI-01",
 domain="Fairness",
 category="Assessment ",
 question=" System Fairness Assessment ?",
 criteria="Fairness Assessment Report",
 ),
 AuditCheckItem(
 item_id="FAI-02",
 domain="Fairness",
 category=" ",
 question="4/5 (DP Ratio >= 0.8) ?",
 criteria=" ",
 ),
 AuditCheckItem(
 item_id="FAI-03",
 domain="Fairness",
 category=" ",
 question=" Model ?",
 criteria=" ",
 ),
 AuditCheckItem(
 item_id="FAI-04",
 domain="Fairness",
 category="Monitoring",
 question="Fairness Monitoring ?",
 criteria="Monitoring , ",
 ),
 ])
 
 # 6. 
 items.extend([
 AuditCheckItem(
 item_id="XAI-01",
 domain=" ",
 category="Implementation",
 question=" System XAI Implementation ?",
 criteria="SHAP/LIME Implementation ",
 ),
 AuditCheckItem(
 item_id="XAI-02",
 domain=" ",
 category=" ",
 question=" ?",
 criteria=" ",
 ),
 AuditCheckItem(
 item_id="XAI-03",
 domain=" ",
 category=" ",
 question=" ?",
 criteria=" , ",
 ),
 ])
 
 # 7. items.extend([
 AuditCheckItem(
 item_id="HUM-01",
 domain=" ",
 category="Level Definition",
 question=" Level(HITL/HOTL/HIC) Definition ?",
 criteria=" , System Definition",
 ),
 AuditCheckItem(
 item_id="HUM-02",
 domain=" ",
 category=" ",
 question="Definition Level ?",
 criteria=" , ",
 ),
 AuditCheckItem(
 item_id="HUM-03",
 domain=" ",
 category=" ",
 question="AI ?",
 criteria=" ",
 ),
 ])
 
 # 8. Document items.extend([
 AuditCheckItem(
 item_id="DOC-01",
 domain="Document ",
 category=" ",
 question=" Document ?",
 criteria="Document Checklist",
 ),
 AuditCheckItem(
 item_id="DOC-02",
 domain="Document ",
 category=" ",
 question=" 5years Management ?",
 criteria=" System, ",
 ),
 AuditCheckItem(
 item_id="DOC-03",
 domain="Document ",
 category=" ",
 question=" / ?",
 criteria=" ",
 ),
 ])
 
 # 9. items.extend([
 AuditCheckItem(
 item_id="USR-01",
 domain=" ",
 category=" ",
 question="AI ?",
 criteria=" , ",
 ),
 AuditCheckItem(
 item_id="USR-02",
 domain=" ",
 category=" ",
 question=" ?",
 criteria=" , ",
 ),
 AuditCheckItem(
 item_id="USR-03",
 domain=" ",
 category=" ",
 question=" ?",
 criteria=" , Period",
 ),
 ])
 
 # 10. Management
 items.extend([
 AuditCheckItem(
 item_id="INC-01",
 domain=" Management",
 category=" ",
 question=" Response ?",
 criteria="Response Document",
 ),
 AuditCheckItem(
 item_id="INC-02",
 domain=" Management",
 category=" ",
 question=" .Management ?",
 criteria=" ",
 ),
 AuditCheckItem(
 item_id="INC-03",
 domain=" Management",
 category=" ",
 question=" ?",
 criteria=" ",
 ),
 ])
 
 return items
 
 def evaluate_item(
 self,
 item_id: str,
 status: ComplianceStatus,
 evidence: str,
 finding: str = "",
 recommendation: str = "",
 ) -> None:
 """Items Assessment"""
 
 for item in self.items:
 if item.item_id == item_id:
 item.status = status
 item.evidence = evidence
 item.finding = finding
 item.recommendation = recommendation
 break
 
 def get_summary(self) -> Dict:
 """ """
 
 total = len(self.items)
 status_counts = {}
 
 for item in self.items:
 status_counts[item.status.value] = status_counts.get(item.status.value, 0) + 1
 
 # Compliance Calculation ( None, Check )
 applicable = total - status_counts.get(" None", 0) - status_counts.get(" Check", 0)
 compliant = status_counts.get("Compliance", 0)
 
 compliance_rate = compliant / applicable if applicable > 0 else 0
 
 # Analysis
 domain_results = {}
 for item in self.items:
 if item.domain not in domain_results:
 domain_results[item.domain] = {'total': 0, 'compliant': 0}
 domain_results[item.domain]['total'] += 1
 if item.status == ComplianceStatus.COMPLIANT:
 domain_results[item.domain]['compliant'] += 1
 
 return {
 'total_items': total,
 'status_counts': status_counts,
 'compliance_rate': compliance_rate,
 'domain_results': domain_results,
 'findings_count': len([i for i in self.items if i.finding]),
 }
 
 def get_findings(self) -> List[Dict]:
 """ """
 
 findings = []
 
 for item in self.items:
 if item.status in [ComplianceStatus.NON_COMPLIANT, ComplianceStatus.PARTIALLY_COMPLIANT]:
 findings.append({
 'item_id': item.item_id,
 'domain': item.domain,
 'question': item.question,
 'status': item.status.value,
 'finding': item.finding,
 'recommendation': item.recommendation,
 })
 
 return findings
 
 def generate_report(self, audit_info: Dict) -> str:
 """Audit Report Generation"""
 
 summary = self.get_summary()
 findings = self.get_findings()
 
 report = f"""
# AI Governance Audit Report

---

## 1. Audit 

| Items | |
|------|------|
| Audit ID | {audit_info.get('audit_id', '')} |
| Audit | {audit_info.get('audit_name', '')} |
| Audit Period | {audit_info.get('period', '')} |
| Audit | {audit_info.get('lead_auditor', '')} |
| Report days | {audit_info.get('report_date', '')} |

---

## 2. ### AI Governance Compliance {summary['compliance_rate']:.1%} .

| Status | Items |
|------|--------|
"""
 
 for status, count in summary['status_counts'].items():
 report += f"| {status} | {count} |\n"
 
 report += f"""

### Major {len(findings)} .

"""
 
 # Compliance 
 non_compliant = [f for f in findings if f['status'] == ' Compliance']
 partial = [f for f in findings if f['status'] == ' Compliance']
 
 if non_compliant:
 report += f"- Compliance: {len(non_compliant)} ( )\n"
 if partial:
 report += f"- Compliance: {len(partial)} ( )\n"
 
 report += f"""

---

## 3. | | Check Items | Compliance | Compliance |
|------|----------|------|--------|
"""
 
 for domain, result in summary['domain_results'].items():
 rate = result['compliant'] / result['total'] if result['total'] > 0 else 0
 report += f"| {domain} | {result['total']} | {result['compliant']} | {rate:.0%} |\n"
 
 report += f"""

---

## 4. """
 
 for i, finding in enumerate(findings, 1):
 severity = "" if finding['status'] == ' Compliance' else ""
 report += f"""
### 4.{i} [{finding['item_id']}] {finding['domain']}

Check Items: {finding['question']}

Status: {severity} {finding['status']}

 : 
{finding['finding']}

 :
{finding['recommendation']}

---
"""
 
 report += f"""

## 5. ### ( Compliance)

"""
 
 for finding in non_compliant:
 report += f"- [{finding['item_id']}] {finding['recommendation']}\n"
 
 report += f"""

### ( Compliance)

"""
 
 for finding in partial:
 report += f"- [{finding['item_id']}] {finding['recommendation']}\n"
 
 report += f"""

---

## 6. | | ID | | | days |
|------|------------|----------|----------|----------|
"""
 
 for i, finding in enumerate(findings, 1):
 report += f"| {i} | {finding['item_id']} | | | |\n"
 
 report += """

---

## 7. 

1. Audit Checklist 
2. 3. ---

 : Audit 
 : Audit 
Approval: Audit 

* Report 15 5years .*
"""
 
 return report
\end{lstlisting}


%----------------------------------------------------------
\section{외부 감사 및 인증}
\label{sec:11.3}

\subsection{외부 감사 준비}
\label{sec:11.3.1}

외부 감사는 독립적인 제3자가 수행하는 감사로, 
최고 수준의 객관성을 제공한다.

\paragraph{외부 감사 유형}

\begin{table}[htbp]
\centering
\caption{외부 감사 유형 및 특징}
\label{tab:external_audit_types}
\begin{tabularx}{\textwidth}{p{3cm}p{3cm}X}
\toprule
\textbf{유형} & \textbf{수행 Institution} & \textbf{특징} \\
\midrule
인증 감사 & 
인증 기관 & 
ISO/IEC 42001 등 표준 인증을 위한 감사; 
합격/불합격 판정 \\
\midrule
보증 감사 & 
보증 기관 & 
거버넌스 효과성에 대한 독립적 보증; 
보증 Report 발행 \\
\midrule
컴플라이언스 감사 & 
법률/컴플라이언스 전문 기관 & 
특정 규제 준수 여부 검증; 
법적 의견 제공 \\
\midrule
알고리즘 감사\cite{saleiro2018aequitas}\cite{metaxa2021auditing} & 
AI 전문 감사 기관 & 
Model 공정성, 편향 등 기술적 감사; 
알고리즘 감사 Report \\
\bottomrule
\end{tabularx}
\end{table}


\paragraph{외부 감사 준비 체크리스트}

\begin{enumerate}
 \item \textbf{문서 정비}: 
 정책, 절차, 기록 등 모든 문서의 완전성 및 현행화 확인
 
 \item \textbf{사전 자가 Check}: 
 내부 감사 또는 모의 감사를 통해 갭 사전 식별 및 보완
 
 \item \textbf{증적 준비}: 
 각 제어 항목에 대한 증거 자료 수집 및 정리
 
 \item \textbf{담당자 교육}: 
 인터뷰 대상자 사전 교육, 예상 질문 준비
 
 \item \textbf{System 접근}: 
 감사인의 시스템 접근 권한, 데모 환경 준비
 
 \item \textbf{일정 조율}: 
 감사 일정, 참석자, 장소 확정
\end{enumerate}


\subsection{ISO/IEC 42001 인증}
\label{sec:11.3.2}

ISO/IEC 42001:2023은 AI 관리 System(AIMS)에 대한 국제 표준으로, 
AI 거버넌스의 체계적 관리를 위한 프레임워크를 제공한다.

\begin{table}[htbp]
\centering
\caption{ISO/IEC 42001 주요 요구사항}
\label{tab:iso42001_requirements_audit}
\begin{tabularx}{\textwidth}{p{2cm}p{3cm}X}
\toprule
\textbf{조항} & \textbf{제목} & \textbf{주요 내용} \\
\midrule
4 & 
조직 상황 & 
이해관계자 요구사항, AI 관리 System Scope \\
\midrule
5 & 
리더십 & 
경영진 의지, 정책, 역할.책임 \\
\midrule
6 & 
계획 & 
리스크 평가, 목표 Setup, 변경 계획 \\
\midrule
7 & 
Support & 
Resource, 역량, 인식, 커뮤니케이션, 문서화 \\
\midrule
8 & 
운영 & 
AI 시스템 생명주기 관리, 공급망 관리 \\
\midrule
9 & 
성능 평가 & 
모니터링, 내부 감사, 경영 검토 \\
\midrule
10 & 
개선 & 
부적합 관리, 지속적 개선 \\
\bottomrule
\end{tabularx}
\end{table}


\paragraph{인증 취득 프로세스}

\begin{enumerate}
 \item \textbf{갭 분석}: 현재 상태와 표준 요구사항 간 갭 식별 (2-4주)
 \item \textbf{System 구축}: 정책, 절차, 제어 구현 (3-6개월)
 \item \textbf{내부 감사}: 구축된 시스템 자체 검증 (2-4주)
 \item \textbf{경영 검토}: 경영진의 시스템 검토 및 승인
 \item \textbf{1단계 감사}: 문서 검토, 인증 준비 상태 확인 (1-2일)
 \item \textbf{2단계 감사}: 현장 감사, 실제 이행 확인 (3-5일)
 \item \textbf{인증 발급}: 적합 판정 시 인증서 발급
 \item \textbf{사후 관리}: 연간 사후 감사, 3년 주기 갱신 감사
\end{enumerate}


\subsection{알고리즘 감사}
\label{sec:11.3.3}

알고리즘 감사(Algorithmic 감사)는 AI 시스템의 기술적 측면, 
특히 공정성, 편향, 투명성\cite{arrieta2020explainable}을 독립적으로 검증하는 전문 감사이다.

\begin{lstlisting}[language=Python, caption={Algorithmic Audit Framework}, label={lst:algorithmic_audit}]
"""
Algorithmic Audit Framework

Technical Validation for Fairness, Bias, and Transparency
"""

from dataclasses import dataclass, field
from typing import Dict, List, Optional, Tuple
import numpy as np
import pandas as pd
from datetime import datetime

@dataclass
class AlgorithmicAuditConfig:
 """Algorithm Audit Configuration"""
 
 # Fairness threshold
 dp_ratio_threshold: float = 0.8
 eo_diff_threshold: float = 0.1
 
 # Sample size
 min_sample_size: int = 1000
 min_subgroup_size: int = 100
 
 # Confidence interval
 confidence_level: float = 0.95


@dataclass
class AuditResult:
 """Audit Result"""
 metric_name: str
 value: float
 threshold: float
 passed: bool
 confidence_interval: Tuple[float, float]
 sample_size: int
 details: Dict = field(default_factory=dict)


class AlgorithmicAuditor:
 """
 Algorithmic Auditor
 Conducts independent technical audits
 """
 
 def __init__(self, config: Optional[AlgorithmicAuditConfig] = None):
 self.config = config or AlgorithmicAuditConfig()
 self.results: List[AuditResult] = []
 
 def audit_fairness(
 self,
 y_true: np.ndarray,
 y_pred: np.ndarray,
 sensitive_features: pd.Series,
 ) -> List[AuditResult]:
 """Fairness Audit"""
 
 results = []
 
 # 1. Demographic Parity
 groups = sensitive_features.unique()
 selection_rates = {}
 
 for group in groups:
 mask = sensitive_features == group
 n = mask.sum()
 if n >= self.config.min_subgroup_size:
 rate = y_pred[mask].mean()
 selection_rates[group] = {'rate': rate, 'n': n}
 
 if len(selection_rates) >= 2:
 rates = [v['rate'] for v in selection_rates.values()]
 dp_ratio = min(rates) / max(rates) if max(rates) > 0 else 1.0
 
 # Calculate confidence interval (bootstrap)
 ci = self._bootstrap_ci(
 y_pred, sensitive_features,
 lambda yp, sf: self._calculate_dp_ratio(yp, sf)
 )
 
 results.append(AuditResult(
 metric_name="Demographic Parity Ratio",
 value=dp_ratio,
 threshold=self.config.dp_ratio_threshold,
 passed=dp_ratio >= self.config.dp_ratio_threshold,
 confidence_interval=ci,
 sample_size=len(y_pred),
 details={'group_rates': selection_rates},
 ))
 
 # 2. Equalized Odds
 tpr_by_group = {}
 fpr_by_group = {}
 
 for group in groups:
 mask = sensitive_features == group
 y_t = y_true[mask]
 y_p = y_pred[mask]
 
 # TPR
 pos_mask = y_t == 1
 if pos_mask.sum() >= self.config.min_subgroup_size:
 tpr = y_p[pos_mask].mean()
 tpr_by_group[group] = tpr
 
 # FPR
 neg_mask = y_t == 0
 if neg_mask.sum() >= self.config.min_subgroup_size:
 fpr = y_p[neg_mask].mean()
 fpr_by_group[group] = fpr
 
 if len(tpr_by_group) >= 2:
 tpr_diff = max(tpr_by_group.values()) - min(tpr_by_group.values())
 fpr_diff = max(fpr_by_group.values()) - min(fpr_by_group.values())
 eo_diff = max(tpr_diff, fpr_diff)
 
 results.append(AuditResult(
 metric_name="Equalized Odds Difference",
 value=eo_diff,
 threshold=self.config.eo_diff_threshold,
 passed=eo_diff <= self.config.eo_diff_threshold,
 confidence_interval=(0, 0), # Simplified
 sample_size=len(y_pred),
 details={'tpr_by_group': tpr_by_group, 'fpr_by_group': fpr_by_group},
 ))
 
 self.results.extend(results)
 return results
 
 def _calculate_dp_ratio(self, y_pred, sensitive_features):
 """Calculate DP Ratio"""
 groups = sensitive_features.unique()
 rates = []
 for group in groups:
 mask = sensitive_features == group
 if mask.sum() >= 30:
 rates.append(y_pred[mask].mean())
 return min(rates) / max(rates) if rates and max(rates) > 0 else 1.0
 
 def _bootstrap_ci(self, y_pred, sensitive_features, metric_fn, n_bootstrap=1000):
 """Bootstrap Confidence Interval"""
 
 n = len(y_pred)
 bootstrap_values = []
 
 for _ in range(n_bootstrap):
 indices = np.random.choice(n, size=n, replace=True)
 value = metric_fn(y_pred[indices], sensitive_features.iloc[indices])
 bootstrap_values.append(value)
 
 alpha = 1 - self.config.confidence_level
 lower = np.percentile(bootstrap_values, alpha/2 * 100)
 upper = np.percentile(bootstrap_values, (1 - alpha/2) * 100)
 
 return (lower, upper)
 
 def audit_data_representativeness(
 self,
 training_data: pd.DataFrame,
 production_data: pd.DataFrame,
 sensitive_column: str,
 ) -> AuditResult:
 """Data Representativeness Audit"""
 
 train_dist = training_data[sensitive_column].value_counts(normalize=True)
 prod_dist = production_data[sensitive_column].value_counts(normalize=True)
 
 # Calculate distribution difference (Total Variation Distance)
 all_groups = set(train_dist.index) | set(prod_dist.index)
 tvd = 0.5 * sum(
 abs(train_dist.get(g, 0) - prod_dist.get(g, 0))
 for g in all_groups
 )
 
 result = AuditResult(
 metric_name="Data Representativeness (TVD)",
 value=tvd,
 threshold=0.1, # 10% difference allowed
 passed=tvd <= 0.1,
 confidence_interval=(0, 0),
 sample_size=len(training_data) + len(production_data),
 details={
 'training_distribution': train_dist.to_dict(),
 'production_distribution': prod_dist.to_dict(),
 },
 )
 
 self.results.append(result)
 return result
 
 def audit_model_stability(
 self,
 predictions_t1: np.ndarray,
 predictions_t2: np.ndarray,
 ) -> AuditResult:
 """Model Stability Audit"""
 
 # Prediction change rate
 changes = predictions_t1 != predictions_t2
 change_rate = changes.mean()
 
 result = AuditResult(
 metric_name="Prediction Stability",
 value=1 - change_rate,
 threshold=0.95, # 95%+ consistency
 passed=(1 - change_rate) >= 0.95,
 confidence_interval=(0, 0),
 sample_size=len(predictions_t1),
 details={'change_rate': change_rate},
 )
 
 self.results.append(result)
 return result
 
 def generate_audit_report(
 self,
 system_info: Dict,
 auditor_info: Dict,
 ) -> str:
 """Audit Report Generation"""
 
 passed_count = sum(1 for r in self.results if r.passed)
 total_count = len(self.results)
 
 report = f"""
# Algorithm Audit Report

---

## 1. Audit Overview

### 1.1 Target System

| Item | Content |
|------|------|
| System Name | {system_info.get('system_name', '')} |
| System ID | {system_info.get('system_id', '')} |
| Version | {system_info.get('version', '')} |
| Purpose | {system_info.get('purpose', '')} |

### 1.2 Audit Information

| Item | Content |
|------|------|
| Organization | {auditor_info.get('organization', '')} |
| Auditor | {auditor_info.get('auditor_name', '')} |
| Date | {auditor_info.get('audit_date', '')} |
| Criteria | {auditor_info.get('criteria', '')} |

---

## 2. ### Summary Results

- Total Checkpoints: {total_count}
- Passed: {passed_count}
- Failed: {total_count - passed_count}
- Pass Rate: {passed_count/total_count*100 if total_count > 0 else 0:.1f}%

### """
 
 if passed_count == total_count:
 report += "[PASS] All fairness criteria met.\n"
 elif passed_count / total_count >= 0.8:
 report += "[WARN] Most criteria met, but some improvements required.\n"
 else:
 report += "[FAIL] Multiple fairness criteria failed. Immediate action required.\n"
 
 report += f"""

---

## 3. Detailed Results

| Metric | Measured Value | Threshold | Result | Conf. Interval |
|--------|--------|------|------|----------|
"""
 
 for result in self.results:
 status = "[PASS] Pass" if result.passed else "[FAIL] Fail"
 ci = f"[{result.confidence_interval[0]:.3f}, {result.confidence_interval[1]:.3f}]" if result.confidence_interval != (0, 0) else "-"
 
 threshold_str = f">= {result.threshold}" if "Ratio" in result.metric_name or "Stability" in result.metric_name else f"<= {result.threshold}"
 
 report += f"| {result.metric_name} | {result.value:.4f} | {threshold_str} | {status} | {ci} |\n"
 
 report += f"""

---

## 4. Detailed Analysis

"""
 
 for i, result in enumerate(self.results, 1):
 report += f"""
### 4.{i} {result.metric_name}

Measured Value: {result.value:.4f} 
Threshold: {'>=' if 'Ratio' in result.metric_name or 'Stability' in result.metric_name else '<='} {result.threshold} 
Result: {'Pass' if result.passed else 'Fail'} 
Sample Size: {result.sample_size:,}

"""
 
 if 'group_rates' in result.details:
 report += "Group Selection Rates:\n\n"
 report += "| Group | Rate | Sample Size |\n|------|--------|--------|\n"
 for group, stats in result.details['group_rates'].items():
 report += f"| {group} | {stats['rate']:.4f} | {stats['n']:,} |\n"
 report += "\n"
 
 report += f"""
---

## 5. Recommendations

"""
 
 failed_results = [r for r in self.results if not r.passed]
 
 if failed_results:
 for result in failed_results:
 report += f"- {result.metric_name}: Threshold not met. Re-training or bias mitigation recommended.\n"
 else:
 report += "- Currently within thresholds; continuous monitoring recommended.\n"
 
 report += f"""

---

## 6. Disclaimer

This audit was performed based on provided data and information.
The results reflect the state at the time of the audit and may vary
with subsequent system or data changes.

---

Organization: {auditor_info.get('organization', '')} 
Auditor: {auditor_info.get('auditor_name', '')} 
Date: {auditor_info.get('audit_date', '')}

*This report is retained for 5 years per Article 15.*
"""
 
 return report
\end{lstlisting}


%----------------------------------------------------------
\section{감사 Response 및 후속 관리}
\label{sec:11.4}

\subsection{지적사항 관리}
\label{sec:11.4.1}

\begin{table}[htbp]
\centering
\caption{지적사항 등급 분류}
\label{tab:finding_classification}
\begin{tabularx}{\textwidth}{p{2cm}p{3cm}X}
\toprule
\textbf{등급} & \textbf{Definition} & \textbf{Response 기한} \\
\midrule
Critical & 
즉각적인 리스크; 규제 위반; 
중대한 제어 실패 & 
즉시 (1주 이내) \\
\midrule
High & 
중요 제어 미비; 
잠재적 규제 위반 & 
단기 (1개월 이내) \\
\midrule
Medium & 
제어 개선 필요; 
효과성 저하 & 
중기 (3개월 이내) \\
\midrule
Low & 
경미한 개선 기회; 
모범 사례 미적용 & 
장기 (6개월 이내) \\
\midrule
Observation & 
참고 사항; 
권고 수준 & 
권고 (선택적) \\
\bottomrule
\end{tabularx}
\end{table}


\subsection{개선 계획 수립 및 이행}
\label{sec:11.4.2}

\begin{table}[htbp]
\centering
\caption{[템플릿 11-1] 감사 지적사항 개선 계획}
\label{tab:template_remediation_plan}
\begin{tabularx}{\textwidth}{|p{3cm}|X|}
\hline
\multicolumn{2}{|c|}{\textbf{감사 지적사항 개선 계획}} \\
\hline
지적사항 ID & \\
\hline
지적사항 제목 & \\
\hline
지적사항 등급 & $\square$ Critical $\square$ High $\square$ Medium $\square$ Low \\
\hline
지적사항 내용 & \\
\hline
근본 원인 & \\
\hline
개선 조치 & \\
\hline
담당 부서 & \\
\hline
담당자 & \\
\hline
목표 완료일 & \\
\hline
필요 Resource & \\
\hline
검증 방법 & \\
\hline
현재 상태 & $\square$ 미착수 $\square$ 진행 중 $\square$ 완료 $\square$ 지연 \\
\hline
완료일 & \\
\hline
검증 결과 & \\
\hline
\end{tabularx}
\end{table}


\subsection{지속적 감사 체계}
\label{sec:11.4.3}

일회성 감사가 아닌 지속적인 감사 체계를 구축한다.

\begin{table}[htbp]
\centering
\caption{지속적 감사 체계}
\label{tab:continuous_audit}
\begin{tabularx}{\textwidth}{p{3cm}p{2.5cm}X}
\toprule
\textbf{활동} & \textbf{주기} & \textbf{내용} \\
\midrule
자동화 모니터링 & 
실시간/일간 & 
성능, 공정성 메트릭 자동 수집 및 임계값 알림 \\
\midrule
자가 Check & 
분기별 & 
시스템 담당자의 체크리스트 based 자가 Check \\
\midrule
샘플 감사 & 
반기별 & 
무작위 선정된 System/결정에 대한 상세 검토 \\
\midrule
내부 감사 & 
연간 & 
독립적인 내부 감사팀의 체계적 감사 \\
\midrule
외부 감사 & 
2년 주기 & 
독립 기관의 인증 또는 보증 감사 \\
\midrule
규제 Response & 
수시 & 
규제 기관 요청 시 즉시 Response \\
\bottomrule
\end{tabularx}
\end{table}


%----------------------------------------------------------
\subsection*{제11장 요약}

본 장에서는 AI 거버넌스 감사 및 인증 체계를 다루었다. 
핵심 내용을 요약하면 다음과 같다.

\begin{enumerate}
 \item \textbf{감사의 가치}: 
 감사는 거버넌스 실효성을 검증하고, 
 개선점을 발굴하며, 이해관계자 신뢰를 제공하는 핵심 메커니즘이다.
 
 \item \textbf{감사 유형}: 
 자가 Check, 내부 감사, 외부 감사, 규제 감사 등 
 다양한 유형이 있으며, 조합하여 활용한다.
 
 \item \textbf{내부 감사}: 
 계획 수립, 현장 감사, 분석 및 평가, 보고의 
 체계적인 방법론을 따른다.
 
 \item \textbf{외부 감사 및 인증}: 
 ISO/IEC 42001 인증, 알고리즘 감사 등을 통해 
 독립적인 검증을 받을 수 있다.
 
 \item \textbf{후속 관리}: 
 지적사항의 등급화, 개선 계획 수립, 이행 추적, 
 지속적 감사 체계 구축이 필요하다.
\end{enumerate}

효과적인 감사 체계는 AI 거버넌스의 지속적인 개선과 
규제 준수를 보장하는 필수 요소이다.
